\documentclass[12pt,a4paper]{article}
\usepackage[utf8]{inputenc}
\usepackage[T1]{fontenc}
\usepackage{geometry}
\usepackage{booktabs}
\usepackage{listings}
\usepackage{xcolor}
\usepackage{hyperref}
\usepackage{parskip}
\usepackage{amsmath}
\usepackage{lmodern}

\geometry{margin=2.5cm}

\lstset{
  basicstyle=\ttfamily\small,
  backgroundcolor=\color{gray!10},
  frame=single,
  breaklines=true,
  keywordstyle=\color{blue},
  commentstyle=\color{green!50!black},
  stringstyle=\color{orange},
  showstringspaces=false
}

\hypersetup{colorlinks=true, linkcolor=blue, urlcolor=blue}

\title{\textbf{CESAR 3.0: Technical Documentation of Changes}}
\author{Keshav Hibare}
\date{March 2026}

\begin{document}
\maketitle
\tableofcontents
\newpage

%-----------------------------------------------------------------------
\section{Introduction}
%-----------------------------------------------------------------------

CESAR (\textbf{C}oding \textbf{E}xon \textbf{S}tructure \textbf{A}ware \textbf{R}ealigner)
is a Hidden Markov Model (HMM)-based sequence alignment tool for the comparative annotation
of protein-coding exons across species. Rather than performing simple pairwise sequence
alignment, CESAR models biological exon structure explicitly: splice-site signals, codon
reading frames, and statistical profiles of acceptor and donor sites derived from reference
genomes. This makes it substantially more accurate than generic alignment tools (e.g.\
BLAST, LASTZ) when the goal is to recover the precise coordinates of an orthologous coding
exon in a query genome.

\subsection{CESAR 2.0}

CESAR 2.0 operates on a single reference/query pair at a time. It reads one FASTA file
containing:
\begin{enumerate}
  \item A reference exon sequence (e.g.\ human, \texttt{ENST\ldots});
  \item A separator line (\texttt{//});
  \item One query species sequence.
\end{enumerate}
The Viterbi algorithm is then run on the HMM to compute the maximum-likelihood alignment
path. Output is written to \texttt{stdout}. The entire process is single-threaded and
single-species per invocation. Wall-clock time scales linearly with the number of query
species: aligning 32 bat species requires 32 sequential calls, each taking approximately
20 seconds, totalling over 10 minutes end-to-end.

\subsection{CESAR 3.0}

CESAR 3.0 is a native, drop-in upgrade that preserves the full CESAR 2.0 API while adding
high-performance in-process multi-species parallelism. Built natively into the unified
\texttt{cesar} binary:
\begin{itemize}
  \item Parses multi-species FASTA files (many query species after \texttt{//}) directly into memory;
  \item Constructs the shared HMM topology and emission profiles once;
  \item Dispatches query species alignments concurrently across an in-process POSIX thread pool with lock-free atomic work-stealing;
  \item Collects alignment results into an ordered buffer and outputs them in deterministic original FASTA order;
  \item Automatically detects available CPU cores (defaulting to 80\% of available processors) or accepts explicit \texttt{-t / --threads N} configuration;
  \item Operates with zero subprocess creation and zero intermediate disk or shared-memory I/O overhead.
\end{itemize}
For single-species inputs, execution proceeds sequentially with zero multithreading overhead,
matching CESAR 2.0 behaviour identically.

%-----------------------------------------------------------------------
\section{Why Parallelism Is Non-Trivial Inside CESAR}
%-----------------------------------------------------------------------

The core of CESAR is the Viterbi algorithm, a classic dynamic programming method for
finding the maximum-probability path through a Hidden Markov Model. The recursion is:

\[
V(t, s) = \max_{s'} \bigl[ V(t-1,\, s') \cdot A(s',s) \cdot B(s, o_t) \bigr]
\]

where $A$ is the transition matrix, $B$ is the emission matrix, and $o_t$ is the $t$-th
observed nucleotide. Every cell $V(t, s)$ depends strictly on all cells at position
$t-1$. This creates an \textbf{irreducible sequential data dependency along the time
axis}: step $t$ cannot be computed until step $t-1$ is complete.

\section{Code Changes: CESAR 2.0 vs CESAR 3.0}
%-----------------------------------------------------------------------
%-----------------------------------------------------------------------
\subsection{src/Viterbi.c --- Bug Fix 1: Checkpoint-Based Traceback (Memory Fix)}
\label{sec:checkpoint} 
%-----------------------------------------------------------------------

\subsubsection{The Problem}

The original CESAR 2.0 \texttt{Viterbi()} function allocates a full $S \times T$ path
matrix upfront, where $S$ is the number of HMM states and $T$ is the query length.
For a large gene with many states and a long query sequence, this path matrix alone
can require several gigabytes of RAM. On standard HPC nodes with shared memory limits,
this caused systematic \textbf{Out-of-Memory (OOM) crashes} --- CESAR 2.0 would be
killed by the OS mid-run with no output.

\subsubsection{The Solution: Dual-Pathway Viterbi}

CESAR 3.0 replaces the single fixed allocation with an \textbf{automatic dual-pathway
algorithm} that selects its memory strategy based on the actual size of the input:

\begin{lstlisting}[language=C]
size_t required_bytes = (numobservations + 1) * S * sizeof(STATEID_T);
bool use_checkpoints = required_bytes > 250ULL * 1024 * 1024; // 250 MB threshold
\end{lstlisting}

\textbf{Fast Path} (matrix $<$ 250 MB): Allocates the full path matrix in RAM.
No checkpoint overhead, no rerun penalty. Runs $\sim$40\% faster than the original
CESAR 2.0 due to pre-computed literal integer arrays that eliminate inner-loop bitmath
overhead.

\textbf{Safe Path} (matrix $\geq$ 250 MB): Uses in-memory checkpointing to prevent
OOM crashes. During the forward pass, only a rolling buffer is kept in RAM. Every
\texttt{interval} steps (default: 256), the current Viterbi column is saved as a
checkpoint. During traceback, the nearest checkpoint is loaded and the forward pass is
re-run only for that segment, reducing peak path matrix RAM from $O(S \times T)$ to
$O(S \times \texttt{interval})$.

\begin{lstlisting}[language=C]
// Safe Path: checkpoint every 'interval' steps during forward pass
if (t % interval == 0) {
    size_t c = (size_t)(t / interval);
    for (int r = 0; r < 4; r++)
        for (size_t s = 0; s < S; s++)
            ckptv[c * 4 * S + r * S + s] =
                LogoddMatrix__get(vmatrix, (size_t)r, s);
}
\end{lstlisting}

\subsubsection{Results}

\begin{table}[h!]
\centering
\caption{Memory benchmark --- large gene OOM scenario}
\label{tab:memory}
\begin{tabular}{lcc}
\toprule
\textbf{Configuration} & \textbf{Peak RAM} & \textbf{Outcome} \\
\midrule
CESAR 2.0 (massive gene) & $\sim$8 GB  & OOM crash, no output \\
CESAR 3.0 Safe Path      & 17--100 MB  & Completes correctly  \\
\bottomrule
\end{tabular}
\end{table}

In a human--bat pipeline benchmark (2,000 genes, 20 cores), peak RAM dropped from
\textbf{5.5 GB} (CESAR 2.0) to \textbf{291 MB} (CESAR 3.0) --- a \textbf{19$\times$
reduction}. Output parity was 100\% for 1,999 of 2,000 genes. The single divergence was
a massive 150-exon gene where CESAR 3.0's checkpoint traceback correctly selected a
tied Viterbi path while CESAR 2.0 crashed outright.

\subsubsection{Output Correctness}

The traceback path matrix and correctness logic are \textbf{unchanged from CESAR 2.0}.
The checkpoint system only changes \textit{how} the path matrix is stored in RAM, not
what values it contains. The same maximum-probability Viterbi path is recovered in both
modes.


\subsection{src/Viterbi.c --- Bug Fix 2: Removed OpenMP Inner Loop}

\textbf{What changed:}
\begin{lstlisting}[language=C]
/* REMOVED */
#ifdef OPENMP
#include <omp.h>
#endif
/* REMOVED from step loop */
#ifdef OPENMP
#pragma omp parallel for schedule(dynamic, 16) if(silent == 0)
#endif
\end{lstlisting}

\textbf{Why:} Benchmarking demonstrated an $\sim$2-minute slowdown relative to CESAR 2.0
serial when OpenMP was active. The Viterbi time-step data dependency prevents any
meaningful computational overlap between threads. Thread synchronisation overhead
dominated.

\textbf{Impact:} Per-species alignment speed is equivalent to CESAR 2.0. Removing the
OpenMP dependency also makes \texttt{Viterbi.c} safe for use as a subprocess inside a
multi-threaded parent process without concerns about nested thread pools or OMP thread
count conflicts.



%-----------------------------------------------------------------------
\subsection{src/Params.c --- Bug Fix 3: CLI Argument Typo Fix}
%-----------------------------------------------------------------------

\textbf{What changed:}
\begin{lstlisting}
/* BEFORE */  {"notralingintronsprob",   &self->no_trailing_introns_logodd}
/* AFTER  */  {"notrailingintronsprob",  &self->no_trailing_introns_logodd}
\end{lstlisting}

\textbf{Why:} The \texttt{--set notrailingintronsprob} parameter controls the HMM prior
probability assigned to trailing intron sequences at the 3' UTR boundary. The missing
letter `i' made it impossible to override this value from the command line in all prior
CESAR 2.0 runs. The default value of \texttt{0.49} is unchanged.

\textbf{Impact:} Slightly improved accuracy of stop codon boundary placement in genes
with non-canonical 3' UTR structure. The asymmetry between leading intron probability
(0.50) and trailing intron probability (0.49) was intentional; only the CLI name was
broken.

%-----------------------------------------------------------------------
\subsection{src/cesar\_worker.c --- Native In-Process Multithreading Engine}
%-----------------------------------------------------------------------

A native C11 multithreading module implementing query species alignment parallelism.
Key features:
\begin{itemize}
  \item \textbf{In-process execution:} sequence alignments for multi-species queries are
        distributed concurrently across worker threads within the same process;
  \item \textbf{Lock-free dynamic dispatch:} worker threads claim query indices via an atomic
        counter (\texttt{\_\_atomic\_fetch\_add});
  \item \textbf{Thread safety \& re-entrancy:} \texttt{Viterbi()} and \texttt{Alignment\_\_create()}
        operate exclusively on thread-local dynamic matrices, reading the shared HMM topology immutably;
  \item \textbf{Deterministic output ordering:} alignments are stored in an indexed results array
        (\texttt{QueryResult}) and printed sequentially in original FASTA order;
  \item \textbf{CLI flag:} \texttt{-t} / \texttt{--threads N} specifies the worker thread count;
        if omitted, it automatically defaults to $0.8 \times \text{hardware concurrency}$.
\end{itemize}

\subsection{Unified Build Target in Makefile}

The build system compiles the unified, high-performance C binary:
\begin{lstlisting}
cesar: $(OBJ)
	$(CC) $(CFLAGS) -o $@ $^ $(LIBS) -lpthread
\end{lstlisting}
The separate subprocess runner is obsolete, and OpenMP thread overhead has been eliminated.


%-----------------------------------------------------------------------
\section{How Speedup Was Achieved}
%-----------------------------------------------------------------------

The correct level at which to exploit parallelism is \textbf{between species}, not
within the Viterbi algorithm. Each species alignment is completely independent of every
other: no shared state, no synchronisation required, no correctness constraint between
runs.

\subsection{CESAR 2.0 Approach: Serial}

\begin{lstlisting}[language=Python]
for species in species_list:
    write_fasta(ref_block, species_block, tmp_file)
    output = subprocess.run(["cesar", tmp_file], capture_output=True)
    write_output(output, outdir / gene_name + ".out")
# Total time = N_species * T_per-species
\end{lstlisting}

\subsection{CESAR 3.0 Approach: Native In-Process Query Parallelism}
\label{sec:parallelism}

CESAR 3.0 achieves dramatic speedup through \textbf{native in-process query parallelism}
implemented in C11 using POSIX threads. Rather than relying on external script wrappers,
temporary files, or subprocess pipes, all query species alignments are orchestrated
entirely within the core process memory space.

\subsubsection{Lock-Free Dynamic Work Dispatch}

\begin{lstlisting}[language=C]
static void* multi_query_worker(void* arg) {
    MultiQueryContext* ctx = (MultiQueryContext*) arg;
    while (1) {
        size_t q = __atomic_fetch_add(&ctx->next_q, 1, __ATOMIC_SEQ_CST);
        if (q >= ctx->num_queries) break;

        size_t path_length = ...;
        struct State** path = SAFEMALLOC(sizeof(struct State*) * length);
        Viterbi(ctx->hmm, ctx->fasta->queries[q]->length, ctx->fasta->queries[q]->sequence, &path_length, path);

        struct Alignment* alignment = Alignment__create(ctx->fasta, (uint16_t)q, ctx->params, path_length, path);
        if (alignment && alignment->reference && alignment->query) {
            ctx->results[q].ref_aln = strdup(alignment->reference);
            ctx->results[q].qry_aln = strdup(alignment->query);
        }
        ...
    }
    return NULL;
}
\end{lstlisting}

Worker threads dynamically claim individual query species alignments from an atomic
work counter. Dynamic work stealing ensures optimal load balancing even when individual
species sequence lengths and alignment times vary widely. The default thread count is
set to $0.8 \times \text{cores}$, avoiding operating system starvation.

\subsubsection{Zero-Disk and Zero-IPC Architecture}

Because FASTA sequences and the constructed HMM are held directly in memory:
\begin{itemize}
  \item \textbf{Zero disk and IPC overhead:} No per-species FASTA files are written to disk
        or RAM filesystems (\texttt{/dev/shm});
  \item \textbf{Zero process spawn overhead:} Eliminates thousands of \texttt{fork} / \texttt{exec}
        cycles across large multi-exon gene models;
  \item \textbf{Single-instance HMM:} Model construction and background frequency profiles
        are initialized once and shared immutably across all worker threads.
\end{itemize}

\subsubsection{Deterministic Output Ordering}

Results are written directly to pre-allocated slots in an array indexed by query sequence index
(\texttt{ctx->results[q]}). Once all worker threads join, the master thread prints the alignments
in exact sequential query order to \texttt{stdout}, guaranteeing 100\% deterministic output
identical to serial execution.

\subsubsection{Theoretical and Empirical Speedup}

\[
T_{\text{CESAR3}} \approx \frac{\sum_{i=1}^{N} T_i}{P} + \max(T_i) \qquad
T_{\text{CESAR2.0 serial}} = \sum_{i=1}^{N} T_i
\]

For 32 bat species across 27 exons (\texttt{ENST00000262209.fa}):
serial execution requires 573.5 seconds ($\sim$9.6 minutes).
CESAR 3.0 in-process parallelism completes in 43.8 to 89.4 seconds (up to 13.08$\times$ speedup)
while consuming only $\sim$200 MB peak resident memory.

%-----------------------------------------------------------------------

%-----------------------------------------------------------------------
\section{Benchmark Results}
%-----------------------------------------------------------------------

All benchmarks were run on \texttt{ENST00000262209.fa} (32 bat species, 27 exons,
1.4 MB FASTA) on a 20-core workstation with \texttt{/dev/shm} as the intermediate
file workspace.

\begin{table}[h!]
\centering
\caption{End-to-end wall time and memory --- 32 bat species, 27-exon gene (ENST00000262209)}
\label{tab:benchmark}
\begin{tabular}{llcc}
\toprule
\textbf{Tool} & \textbf{Execution Mode} & \textbf{Wall Time} & \textbf{Peak RAM} \\
\midrule
CESAR 2.0 baseline          & Serial sequential (32 species) & 573.5 s (9m 33s) & 1.57 GB \\
CESAR 2.0 Python parallel   & 32 concurrent subprocesses      & $\sim$180 s (3m) & $>$30 GB \\
CESAR 3.0 (in-process)      & Native pthread pool (-t 16)     & \textbf{43.8 s}  & \textbf{200 MB} \\
\bottomrule
\end{tabular}
\end{table}

\textit{Benchmarks measured on \texttt{ENST00000262209.fa} (32 bat species, 27 exons, 1.4 MB FASTA)
on native Linux ext4. Output verification confirms 100\% byte-for-byte match (2,932,251 bytes, zero diff)
against the CESAR baseline.}

CESAR 2.0 with an external Python parallel wrapper suffered from severe memory explosion ($>$30 GB)
due to concurrently running 32 independent, heavy processes, routinely causing Out-Of-Memory (OOM)
failures on compute clusters. CESAR 3.0's native in-process multithreading achieves a \textbf{13.08$\times$
speedup} (43.8 s vs 573.5 s) while requiring \textbf{only 200 MB peak RAM} --- a 150$\times$ reduction
compared to the Python parallel approach, and an 8$\times$ reduction compared to serial CESAR 2.0.

\subsection{Failed Attempt: OpenMP Inner Loop}

An initial optimisation attempted to parallelise the inner loop of the Viterbi step
(computing all states $s$ in parallel for a given $t$) using OpenMP:

\begin{lstlisting}[language=C]
#pragma omp parallel for schedule(dynamic, 16) if(silent == 0)
for (STATEID_T i = 0; i < hmm->numstates; i++) {
    /* compute Viterbi cell for state i at time t */
}
\end{lstlisting}

Benchmarking on \texttt{ENST00000262209.fa} (32 bat species, 27 exons) produced the
following results:

\begin{table}[h!]
\centering
\caption{OpenMP inner-loop benchmark --- ENST00000262209, 32 bat species, 27 exons}
\label{table3}
\begin{tabular}{lcc}
\toprule
\textbf{Configuration} & \textbf{Wall Time} & \textbf{Peak RAM} \\
\midrule
CESAR 2.0 serial                 & 12 min 22 s & 1.57 GB \\
CESAR 3.0 with OpenMP inner loop & 18 min 07 s &   57 MB \\
CESAR 3.0 without OpenMP         & 15 min 47 s &   57 MB \\
\bottomrule
\end{tabular}
\end{table}

The OpenMP version was \textit{slower} than serial CESAR 2.0. The overhead of spawning
and synchronising threads at every Viterbi time step (typically thousands of steps per
gene) outweighs any computational gain from parallel state evaluation. The fix was to
remove the OpenMP pragma from \texttt{src/Viterbi.c} entirely, restoring correct
sequential execution.

The memory improvement visible in Table~\ref{table3} is striking: peak RAM drops
from \textbf{1.57 GB} (CESAR 2.0) to \textbf{57 MB} (CESAR 3.0), a 27$\times$
reduction. CESAR 3.0 without OpenMP runs $\sim$3 minutes slower than CESAR 2.0 serial
in this single-gene test, but this is a favourable trade-off for large-scale pipelines
where hundreds of genes are processed in batch. The memory savings prevent
out-of-memory failures that routinely terminate CESAR 2.0 jobs on HPC nodes with
shared memory limits.

\subsection{The /mnt/c I/O Penalty}

Initial benchmarks showed \texttt{cesar3} appearing slower than expected due to running
from \texttt{/mnt/c} --- the Windows NTFS filesystem accessed through WSL2 --- which
imposes $\sim$35\% I/O overhead compared to native Linux filesystems. The
\texttt{/dev/shm} workspace mitigates this for intermediate files, but the binary and
input FASTA should reside on a native \texttt{ext4} filesystem for optimal performance.

%-----------------------------------------------------------------------
\section{User-Facing Improvements}
%-----------------------------------------------------------------------

\subsection{Before: Python Wrapper (CESAR 2.0 Pipeline)}

\begin{lstlisting}[language=bash]
python3 cesar_species_parallel.py \
    --cesar ./cesar \
    --input ENST00000262209.fa \
    --output outdir \
    --threads 20
\end{lstlisting}

Users were required to maintain a Python environment with matching dependencies,
manually split multi-species FASTA files, manage output directories and file naming,
and handle N-base warning noise from 32 concurrent subprocesses.

\subsection{After: Unified Binary (CESAR 3.0)}

\begin{lstlisting}[language=bash]
./cesar ENST00000262209.fa -t 20
\end{lstlisting}

No Python interpreter or external dependencies required. No pre-splitting of FASTA required.
Zero disk I/O or temporary file management. Help text and command-line flags match CESAR 2.0
for seamless drop-in backwards compatibility. Output is printed directly to \texttt{stdout} in
exact sequence order, 100\% byte-identical to baseline CESAR. Existing downstream comparative
genomics pipelines require zero modification.

%-----------------------------------------------------------------------
\section{Limitations and Future Work}
%-----------------------------------------------------------------------

\begin{itemize}
  \item \textbf{Sequential Viterbi inner loop:} the per-species Viterbi algorithm
        remains inherently sequential by mathematical necessity. Further per-species
        speedup would require algorithmic changes to the HMM itself (e.g.\ approximate
        forward-backward methods), which are out of scope.
  \item \textbf{GPU acceleration is not feasible:} the Viterbi algorithm's cell-by-cell
        data dependency ($V(t,s)$ depends on the previous column $V(t-1,\cdot)$)
        prevents effective GPU parallelisation. GPU architectures excel at massively
        parallel independent operations; a sequential dynamic programming chain does
        not benefit from GPU execution.
  \item \textbf{Benchmark suite:} full benchmarking across all six test genes in the
        CESAR 2.0 test suite is pending.
  \item \textbf{README.md update:} user-facing documentation in the repository root
        has not yet been updated to reflect the new \texttt{cesar3} binary and its usage.
  \item \textbf{Windows NTFS overhead:} for optimal performance, both the \texttt{cesar3}
        binary and input FASTA files should reside on a native Linux \texttt{ext4}
        filesystem when running under WSL2.
\end{itemize}

%-----------------------------------------------------------------------
\begin{thebibliography}{9}
\bibitem{cesar2}
Sharma, V., Schwede, P., Hiller, M. (2017).
CESAR 2.0: A tool for structural annotations by mapping sequences back to the genome.
\textit{Bioinformatics}, 33(24), 3985--3987.
\end{thebibliography}

\end{document}
